\section{The Metaphysical Principles of the Infintesimal Calculus}

As someone who awoke rather late to the importance of mathematics (whose general use for the esoteric student is not only helpful, but even necessary), I decided to read my first full Rene Guenon treatise\footnote{\url{https://arcaneknowledgeofthedeep.files.wordpress.com/2014/02/guenonmetaphysicalprinciplesinfinitesimalcalculus.pdf}}.


The fundamental distinction which he makes is between the Indefinite and the Infinite. Since Wilhelm Gottfried Leibniz\footnote{\url{http://en.wikipedia.org/wiki/Gottfried\_Wilhelm\_Leibniz}} was the progenitor of the modern calculus, and since the term ``numerable infinity'' (or rather the concept) originated with him, Guenon takes aim primarily at Leibniz's explanations and way of thinking. Guenon's contention is that all of the criticism (and the subsequent development of the field of calculus) could have been adequately dealt with, had the proper distinction been made between what is Indefinite, and what is truly Infinite. \textbf{Guenon insisted that there was only one Infinite, and it was not divisible or numerable. It follows from this that any Number is necessarily non-Infinite, since whatever is one thing rather than another (as all Number implies about itself) is by definition non-Infinite. It is therefore, numerable or quantifiable, and therefore Indefinite, rather than Infinite.}

Leibniz had a confused notion of esoteric studies (Rosicrucian) and tried to modify his terms to suit the understanding of those he was speaking to, but Guenon contended that he did this rather poorly, as he failed to make the distinction fundamental to any science, and so opened the path for the modern tendency to do away with the one Infinite in favor of quantifiable infinities or multiple infinities on the secular or profane plane of being.

One can see this confusion in the subject of Fractals, which are somehow implied to be infinite (given that they can always be expanded to see ``more'' articulation, which repeats itself ``endlessly''). This endlessness is not infinite, for it belongs to the plane of the finite, and represents an ability of the human mind to divide the finite indefinitely, rather than truly opening up some kind of substitute infinite perspective that can displace the one Infinite (which is impossible).

Guenon deals with the subject in some technical depth, which ought to interest those inclined towards mathematics by nature and training, and is a master (even in translation) of maintaining proper distinctions and speaking accurately: as they once said of Dr Greg Bahnsen\footnote{\url{http://en.wikipedia.org/wiki/Greg_Bahnsen}}, \emph{rem acu tetigisti} (\emph{you touched the thing with a needle}).


Leibniz wrestled with the implications of what he himself was trying to say, but ultimately was lead astray in his thinking by his own terminology. The concept of ``limit'' is actually (according to Guenon) analogous to the transition from one plane (the physical) to that of another (the metaphysical), and represents a ``transposition'' from a limited state of being to that of an infinite one. I believe there are some posts at Gornahoor on that topic, under Michaelstadter. Nevertheless, the notion of ``limit'' when used mathematically is not properly infinite, but connotes the event horizon of what is indefinite, and teaches by analogy what is proper to the metaphysical sphere only.

Number can increase indefinitely, and in fact (as Archimedes showed long ago\footnote{\url{http://en.wikipedia.org/wiki/The\_Sand\_Reckoner}}) can be calculated up to a point at which the human mind can no longer comprehend what it is calculating. That is, it can go on indefinitely, until the human mind is ``mazed'' and dissolves. But since it is in fact a Number that is being reckoned, it is not, and never can be, Infinite.

On a related note, I picked up Phillip Pullman's \emph{The Golden Compass}, the first book in \emph{His Dark Materials}, which is a quote from John Milton (who was ``of the devil's party without knowing it'', according to Blake).

\begin{quotex}
\begin{verse}
Into this wild Abyss\\
The womb of Nature, and perhaps her grave\\
Of neither sea, nor shore, nor air, nor fire,\\
But all these in their pregnant causes mixed\\
Confusedly, and which thus must ever fight,\\
Unless the Almighty Maker them ordain\\
His dark materials to create more worlds,\\
Into this wild Abyss the wary Fiend\\
Stood on the brink of Hell and looked a while,\\
Pondering his voyage; for no narrow frith\\
He had to cross.
\end{verse}
\end{quotex}


Pullman hated CS Lewis, and wrote his ``atheist fantasy'' to displace the Narnian worldview. He is in the interesting intellectual position of refuting himself, since once one has experienced the ``true God'' as a ``Demiurge'', one can thereafter never really be certain that 1. You aren't experiencing an ``alternate'', deeper illusion of the Demi-Urge (his ``Loving Side'') or 2. that there isn't another truer God behind the Light of Lights, which would then turn into a second, higher Demi-Urge. Or, as Tomberg writes, ``there are illusions above, as there are illusions below''.

Pullman is a perfect example of what happens when one sets out to defend a true, but false idea (in this case, self-contradictory) with psychologically plausible fantasy. it is, in fact, true that many people worship a ``false God'' as the true Creator, but this fact is recognized in all Traditional science and in all sacred Scriptures. As good as the novel is (it is very taut and well written, with interesting ideas to play with), it turns out that (additionally) that if Pullman's ideology triumphs in Reality, one will not be able to salvage enough Romance and aura of danger to set further fantasy novels, since the entire ambience of his novel revolves around an alternate reality ``High Church England'', replete with dreamy spires and dread castles, which don't exist anymore once the ``new man'' has created the brave new world of the Revolution. As always, the first circles of Hell are the most romantic and exciting. After that, it becomes a bit too predictable. Readers of Pullman's novels who wish to ``deconstruct'' him will have to imagine, first, The Cathedral\footnote{\url{http://unqualified-reservations.blogspot.co.uk/2009/01/gentle-introduction-to-unqualified.html}} \emph{a la moderne}. Which requires more intellectual or spiritual horsepower than (say) smashing the Middle Ages, since a nose on one's face is harder to see. .

There are, however, indefinite, or countless fantasies, which Revolution can dream up, and only one way to be Right. It takes a genius like Guenon to see, diagnose, and dissect this error as it manifested in the secular science of modern mathematics.


\flrightit{Posted on 2015-05-09 by Logres}

